\section{Aim 3 --- Bluesky steering framework}
\label{sec:aim3}

\textit{Strawman owners: Campbell, with Allan, Maffettone, Caswell, Rakitin,
Barbour, Olds. $\sim$500 words.}

Aim~3 is the shared loop, not a new control system. The four steps in the
proposal are reduce, decide the next action, encode it to controls, and
collect. Two timescales sit inside that loop: faster than a measurement
(instrument optimization) and commensurate with a measurement (sequence
planning). Interoperation was specified with the Queue Server, workflow
systems, and parameterized similar measurements (for example powder). The UI
was specified to show the live experiment, agent state, history, queued plans,
and hyper-parameters \cite{allan2019bluesky}.

Bluesky was the reference implementation \emph{because} pieces can be adopted
without replacing a beamline. This paper should score that interop claim, not
assume a single stack won.

The 2023 ``Agents'' are pluggable modules in this pipeline --- local or
facility/HPC --- that consume reduced data and emit the next measurement or
control directive. Gaussian-process fly-scan demos at NSLS-II already existed
when the proposal was written. They are not language models.

\textit{Placeholders:} Bluesky/Tiled status; LCLS MFX photon-beam tuning and
operator-facing skills as one Aim~3 + Aim~1 crystallography cell.


\subsection{Blop: a short history}
\textit{Owner: Hopkins.}

Blop (the Bluesky Optimization Package) began at NSLS-II as a single-purpose tool:
a custom-built wrapper around BoTorch for Gaussian-process Bayesian optimization,
with little integration into Bluesky itself, for autonomously aligning a beamline
by maximizing a measured figure of merit. That design was published in the
Journal of Synchrotron Radiation in 2024 as a general algorithm for autonomous
beamline alignment \cite{Morris2024}. Starting in 2025, the project outgrew that
single-purpose shape and became a general-purpose steering toolbox for Bluesky.

That work did two things at once: it replaced the hand-built BoTorch wrapper with
Ax, which handles trial generation and experiment tracking on top of BoTorch; and
it separated the optimization logic from the Bluesky-facing code behind a small
set of protocols, so that a plan-facing half that only knows how to run a Bluesky
acquisition and evaluate its outcome stays separate from an optimizer-facing half
that proposes the next point, checkpoints its state, and recovers from a failed
trial. Once that boundary existed, Blop became a bridge between optimizer backends
and Bluesky, via those protocols and the plans built on them, rather than a single
Ax/BoTorch agent: new backends could be added without touching the Bluesky side,
including a classical SciPy backend (local minimization, SHGO, dual annealing) for
problems that do not need a surrogate model, and an Xopt backend, alongside the
Ax/BoTorch Bayesian optimizer.

The same year of work hardened Blop for the kind of unattended, remote operation a
multi-facility steering framework needs: optimizer state can be checkpointed and
restored, a failed evaluation no longer stalls a run, and an optimization can stop
itself against a global criterion instead of a fixed trial count. Its integration
with the Bluesky Queue Server grew from an experimental add-on into a restartable,
checkpointable component that can run on a remote queue unattended. None of this
removed the experimenter from the loop: an operator can still override a suggested
point manually, and review model fit and uncertainty before a move is committed.

Blop also grew as a project rather than just as code: it moved out of a single
facility's repository into the shared \texttt{bluesky} GitHub organization,
reached its first stable 1.0 release in mid-2026, and is now packaged for both PyPI
and conda-forge. What began as one beamline's alignment script is now a shared,
multi-backend building block for the decision-support layer this Aim describes.

The next step builds on that same core rather than replacing it: Blop's
protocol-based architecture is migrating into \texttt{bluesky-adaptive}, enabling
service-oriented deployment and passive, off-policy learning agents alongside the
online optimizers described above. Once that migration lands, the Blop package
itself will be archived; the architecture built over the past year carries forward
unchanged.

\subsubsection{Blop development}

Blop, or Bluesky Optimization Package \cite{Morris2024} is an autonomous agent 
framework designed to steer experimental data acquisition by bridging advanced
optimization algorithms with the Bluesky ecosystem. Built upon the adaptive 
experimentation platform Ax platform \cite{olson2025ax}, Blop provides a 
streamlined interface that simplifies how experimental data is collected, analyzed, 
and formulated as an optimization problem. 

With the release of version 1.0 in July 2026 on conda-forge and PyPI, Blop introduced 
significant architectural and functional enhancements. Key updates include revamped 
documentation featuring active deployments, tutorials, and comprehensive how-to guides. 
Core functionality has been expanded to incorporate first-class integrations with 
Bluesky Queueserver, Xopt, and SciPy. Additionally, simulation capabilities now feature 
a “bring your own beamline” workflow utilizing XRT to accurately model physical beamline 
movements. Advanced optimization features further enhance experimental flexibility, 
introducing scalarized objectives, manual data point insertion and data priming, best 
point navigation, and robust checkpointing. 

\subsubsection{Blop-sim and XRT}

Blop-sim is a currently in development add-on to blop targeting offline tuning and deployment
testing. Deployment of Blop workflows require discretionary beamtime like classical 
experiments to validate their performance, and the less time needed to do online setup the 
more algorithmic tuning and optimizer configuration can be done. There are also deployment
factors which can easily be covered with a rough simulation of the beamline ie:  

\begin{itemize} 

\item Do evaluation functions work and provide the correct cost?  

\item Are we in an order of magnitude of the correct bounds (is an axis too sensitive)?  

\end{itemize} 

This drives us to develop a "Beamline in a Box" system as an aggressively scaled semipermanent 
simulation, which acts as an "as best" software indistinguishable setup allowing for accurate, 
rapid, and reconfigurable offline deployment testing. We target swapout at the epics channel 
layer as it is a convenient flat low-level interface to mock a beamline through, with opaque 
self contained scalars providing an interface for arbitrary simulation complexity.

\subsubsection{Lessons from beamline deployments}
At CHX, Blop has been successfully applied to optimization tasks involving Bragg peaks and 
detector positioning. At XPD, ongoing work demonstrates its potential in autonomous experimental 
workflows. 

An earlier effort at BMM was less successful, highlighting the challenges of moving from an 
optimization concept to a working beamline deployment. Experiences like this reinforce the 
importance of realistic simulations, carefully defined evaluation functions, and appropriate 
parameter bounds before committing valuable beamtime. 

Together, these deployments demonstrate both Blop's capabilities and the practical considerations 
needed for reliable optimization. 